% Generated by publication/build_sections.py; do not edit by hand.
\paragraph{Informal verdict.}
\textbf{A --- Complete solution.} Confidence: \textbf{high (0.94)}.

\paragraph{Lean coverage.}
Full canonical Q855 coverage relative to the premise imported from R773. Lean defines the universal coefficients alphaC by the audited triangular recursion and proves, for every real lambda, every recurrence solution v satisfying v\(n\)/n\textasciicircum{}lambda -> 1, and every N, the stronger remainder v\(n\)-sum\_\{k<=N\} alphaC(lambda,k)n\textasciicircum{}\(lambda-k\)=O\(n^(lambda-N-1)\). The single theorem covers lambda=0. It does not re-prove R773's existence/basis theorem, which the canonical question states as prior input; it is not the permitted finite-cancellation/natural-lambda fallback.

\subsection{Audited informal answer}
Q855 --- complete solution
\begin{enumerate}
\item Verified statement
\end{enumerate}
On printed pages 14--15 of the canonical PDF, Q855 considers, for a fixed parameter \(\lambda\), the real solutions of

\[
u_{n+2}=u_{n+1}+\frac{\lambda}{n+2}u_n
\qquad(n\ge 0).
\tag{1}
\]

The preceding response R773 is quoted as proving that the two-dimensional solution space has a basis \(\{(v_n),(w_n)\}\) such that

\[
v_n\sim n^\lambda,
\qquad
w_n\sim \frac{(-1)^n\lambda^n}{n!\,n^{\lambda+1}}.
\tag{2}
\]

The question asks for a complete asymptotic expansion of the dominant solution \(v\):

\[
v_n=n^\lambda+\alpha_1n^{\lambda-1}+\cdots+
\alpha_Nn^{\lambda-N}+o(n^{\lambda-N})
\]

for every \(N\). Luc Pommeret+2Luc Pommeret+2
The PDF does not explicitly quantify \(\lambda\). I take \(\lambda\in\mathbb R\) fixed. There is one minor endpoint defect in the quoted premise: when \(\lambda=0\), the displayed expression for \(w_n\) is zero for \(n\ge1\), so it cannot literally be an asymptotic equivalent of a nonzero sequence. The requested conclusion at \(\lambda=0\) is nevertheless immediate and will be included.

\begin{enumerate}
\item Final theorem
\end{enumerate}
For \(x\in\mathbb R\) and \(m\in\mathbb N\), write

\[
\binom{x}{m}=\frac{x(x-1)\cdots(x-m+1)}{m!}.
\]

Define \(\alpha_0(\lambda)=1\), and for \(r\ge1\) define recursively

\[
\boxed{
\alpha_r(\lambda)
=
\frac1r
\sum_{k=0}^{r-1}
\alpha_k(\lambda)
\left[
\bigl(2^{\,r+1-k}-1\bigr)
\binom{\lambda-k}{r+1-k}
-\lambda(-2)^{r-k}
\right].
}
\tag{3}
\]

Then, for every fixed \(\lambda\in\mathbb R\), every solution \(v\) of \(1\) satisfying \(v_n\sim n^\lambda\), and every \(N\ge0\),

\[
\boxed{
v_n
=
\sum_{k=0}^{N}\alpha_k(\lambda)n^{\lambda-k}
+
O_{\lambda,N,v}\!\left(n^{\lambda-N-1}\right).
}
\tag{4}
\]

This is one order stronger than requested.
In particular,

\[
v_n\sim n^\lambda\sum_{k\ge0}\frac{\alpha_k(\lambda)}{n^k}
\]

in the Poincaré sense. The coefficients depend only on \(\lambda\), not on the choice of the normalized dominant solution: two such solutions differ by a multiple of \(w\), which is smaller than every power of \(n^{-1}\).
The first coefficients are

\[
\boxed{\alpha_1(\lambda)=\frac{\lambda(3\lambda+1)}2,}
\tag{5}
\]


\[
\boxed{
\alpha_2(\lambda)
=
\frac{\lambda(\lambda-1)
\left(27\lambda^2+5\lambda+2\right)}{24},
}
\tag{6}
\]


\[
\boxed{
\alpha_3(\lambda)
=
\frac{\lambda^2(\lambda-1)(\lambda-2)
(3\lambda-1)(9\lambda-1)}{48},
}
\tag{7}
\]

and

\[
\boxed{
\alpha_4(\lambda)
=
\frac{\lambda(\lambda-1)(\lambda-2)(\lambda-3)}{5760}
\left(
1215\lambda^4-1890\lambda^3+485\lambda^2-18\lambda-8
\right).
}
\tag{8}
\]


\begin{enumerate}
\item Construction of the formal expansion
\end{enumerate}
Introduce the recurrence operator

\[
(L_\lambda u)_n
=
u_{n+2}-u_{n+1}-\frac{\lambda}{n+2}u_n.
\]

We seek an approximate solution of the form

\[
p(n)=n^\lambda A\!\left(\frac1n\right),
\qquad
A(x)=\sum_{k\ge0}a_kx^k,\quad a_0=1.
\]

Put \(x=1/n\). Direct substitution gives the exact identity

\[
n^{-\lambda}(L_\lambda p)_n
=
\Phi_\lambda(A)(x),
\tag{9}
\]

where

\[
\begin{aligned}
\Phi_\lambda(A)(x)
={}&
(1+2x)^\lambda
A\!\left(\frac{x}{1+2x}\right)
-
(1+x)^\lambda
A\!\left(\frac{x}{1+x}\right) \\
&-
\frac{\lambda x}{1+2x}A(x).
\end{aligned}
\tag{10}
\]

Since

\[
(1+cx)^\lambda
A\!\left(\frac{x}{1+cx}\right)
=
\sum_{k\ge0}a_kx^k(1+cx)^{\lambda-k},
\]

we have

\[
\Phi_\lambda(A)(x)
=
\sum_{k\ge0}a_kx^k
\left((1+2x)^{\lambda-k}-(1+x)^{\lambda-k}\right)
-
\frac{\lambda x}{1+2x}\sum_{k\ge0}a_kx^k.
\tag{11}
\]

The constant coefficient vanishes. The coefficient of \(x\) also vanishes automatically:

\[
2\lambda-\lambda-\lambda=0.
\]

For \(r\ge1\), the coefficient of \(x^{r+1}\) is

\[
\begin{aligned}
[x^{r+1}]\Phi_\lambda(A)
={}&
\sum_{k=0}^{r}
a_k
\bigl(2^{r+1-k}-1\bigr)
\binom{\lambda-k}{r+1-k}\\
&-
\lambda\sum_{k=0}^{r}a_k(-2)^{r-k}.
\end{aligned}
\tag{12}
\]

The contribution involving \(a_r\) is

\[
a_r\left((2-1)(\lambda-r)-\lambda\right)=-r\,a_r.
\]

Consequently, once \(a_0,\ldots,a_{r-1}\) have been fixed, there is a unique choice of \(a_r\) that makes the coefficient of \(x^{r+1}\) vanish. Solving for \(a_r\) gives exactly \(3\).
Thus there is a unique formal series

\[
A_\lambda(x)=\sum_{k\ge0}\alpha_k(\lambda)x^k,
\qquad \alpha_0=1,
\tag{13}
\]

such that

\[
\Phi_\lambda(A_\lambda)=0
\]

as a formal power series.
For fixed \(N\), put

\[
A_{\lambda,N}(x)=\sum_{k=0}^{N}\alpha_k(\lambda)x^k,
\qquad
p_N(n)=n^\lambda A_{\lambda,N}(1/n).
\tag{14}
\]

The coefficient of \(x^j\) in \(\Phi_\lambda(A)\) depends only on
\(a_0,\ldots,a_{j-1}\). Hence the truncation retains all cancellations through degree \(N+1\):

\[
\Phi_\lambda(A_{\lambda,N})(x)=O_{\lambda,N}(x^{N+2})
\qquad(x\to0).
\]

Using \(9\),

\[
\boxed{
(L_\lambda p_N)_n
=
O_{\lambda,N}\!\left(n^{\lambda-N-2}\right).
}
\tag{15}
\]

It remains to prove that an approximate solution with such a small residual is close to the actual dominant solution.

\begin{enumerate}
\item Stability lemma for the recurrence
\end{enumerate}
We use only the asymptotic basis \(2\) supplied in the statement.
Lemma
Let \(\lambda\ne0\), let \(V,W\) be solutions of the homogeneous recurrence \(1\) such that

\[
V_n\sim n^\lambda,
\qquad
W_n\sim \frac{(-1)^n\lambda^n}{n!\,n^{\lambda+1}}.
\tag{16}
\]

Let \(s\ge0\). Suppose that, for all sufficiently large \(n\),

\[
y_{n+2}-y_{n+1}-\frac{\lambda}{n+2}y_n=g_n,
\tag{17}
\]

where

\[
g_n=O(n^{\lambda-s-2}),
\qquad
y_n=o(n^\lambda).
\tag{18}
\]

Then

\[
\boxed{y_n=O(n^{\lambda-s-1}).}
\tag{19}
\]

Proof
Define the Casoratian

\[
\Delta_n=V_{n+1}W_n-V_nW_{n+1}.
\]

From \(16\),

\[
\frac{V_{n+1}}{V_n}\longrightarrow1
\]

and

\[
\frac{W_{n+1}}{W_n}
=
-\frac{\lambda}{n+1}
\left(\frac{n}{n+1}\right)^{\lambda+1}(1+o(1))
=
O(n^{-1}).
\]

Therefore

\[
\begin{aligned}
\Delta_{n+1}
&=
V_{n+2}W_{n+1}-V_{n+1}W_{n+2}\\
&=
V_{n+1}W_{n+1}
\left(
\frac{V_{n+2}}{V_{n+1}}
-
\frac{W_{n+2}}{W_{n+1}}
\right)\\
&\sim V_{n+1}W_{n+1}.
\end{aligned}
\tag{20}
\]

For all sufficiently large \(n\), write

\[
\binom{y_{n+1}}{y_n}
=
\begin{pmatrix}
V_{n+1}&W_{n+1}\\
V_n&W_n
\end{pmatrix}
\binom{A_n}{B_n}.
\tag{21}
\]

The vector recurrence associated with \(17\) gives the standard variation-of-constants identities

\[
A_{n+1}-A_n
=
g_n\frac{W_{n+1}}{\Delta_{n+1}},
\tag{22}
\]


\[
B_{n+1}-B_n
=
-g_n\frac{V_{n+1}}{\Delta_{n+1}}.
\tag{23}
\]

By \(20\),

\[
\frac{W_{n+1}}{\Delta_{n+1}}=O(n^{-\lambda}),
\qquad
\frac{V_{n+1}}{\Delta_{n+1}}
=O\!\left(\frac1{|W_{n+1}|}\right).
\tag{24}
\]

Consequently,

\[
A_{n+1}-A_n=O(n^{-s-2}),
\tag{25}
\]

so \(A_n\) converges to a finite limit \(A_\infty\).
We next control the \(B_nW_n\) contribution. From \(23\),

\[
|B_nW_n|
\le |B_{n_0}W_n|
+
C|W_n|
\sum_{j=n_0}^{n-1}
\frac{j^{\lambda-s-2}}{|W_{j+1}|}.
\tag{26}
\]

Set

\[
T_j=\frac{j^{\lambda-s-2}}{|W_{j+1}|}.
\]

Using \(W_{j+1}/W_j=O(j^{-1})\),

\[
\frac{T_{j-1}}{T_j}
=
\left(\frac{j-1}{j}\right)^{\lambda-s-2}
\frac{|W_{j+1}|}{|W_j|}
=
O(j^{-1}).
\tag{27}
\]

Thus, for large \(j\), \(T_{j-1}\le T_j/2\), and hence

\[
\sum_{j=n_0}^{n-1}T_j=O(T_{n-1}).
\]

It follows from \(26\) that

\[
B_nW_n
=
O(n^{\lambda-s-2}).
\tag{28}
\]

Here the fixed term \(B_{n_0}W_n\) is harmless because \(W_n\) is smaller than every negative power of \(n\): this follows directly from the factorial in \(16\).
Since \(V_n\sim n^\lambda\), \(28\) gives

\[
\frac{B_nW_n}{V_n}=O(n^{-s-2})\longrightarrow0.
\tag{29}
\]

On the other hand, \(21\) gives

\[
\frac{y_n}{V_n}
=
A_n+B_n\frac{W_n}{V_n}.
\]

By hypothesis \(y_n=o(n^\lambda)\), so \(y_n/V_n\to0\). Together with \(29\), this forces

\[
A_\infty=0.
\]

Therefore, by \(25\),

\[
A_n
=
-\sum_{j=n}^{\infty}(A_{j+1}-A_j)
=
O(n^{-s-1}).
\]

Finally,

\[
y_n=A_nV_n+B_nW_n
=
O(n^{\lambda-s-1})+O(n^{\lambda-s-2}),
\]

which proves \(19\). \(\square\)

\begin{enumerate}
\item Application of the stability lemma
\end{enumerate}
Fix \(N\ge0\), and define \(p_N\) by \(14\). Set

\[
y_n=v_n-p_N(n).
\]

Since both \(v_n\) and \(p_N(n)\) are equivalent to \(n^\lambda\),

\[
y_n=o(n^\lambda).
\tag{30}
\]

Moreover, \(L_\lambda v=0\), while \(15\) gives

\[
(L_\lambda y)_n
=
-(L_\lambda p_N)_n
=
O(n^{\lambda-N-2}).
\tag{31}
\]

For \(\lambda\ne0\), the stability lemma with \(s=N\) yields

\[
v_n-p_N(n)=O(n^{\lambda-N-1}),
\]

which is exactly \(4\).
For \(\lambda=0\), recurrence \(1\) reduces to

\[
u_{n+2}=u_{n+1}.
\]

Thus every solution is constant from index \(1\) onward. The normalization \(v_n\sim1\) forces

\[
v_n=1\qquad(n\ge1).
\]

Formula \(3\) gives \(\alpha_r(0)=0\) for every \(r\ge1\), so \(4\) also holds at \(\lambda=0\).
This completes the proof for every real \(\lambda\).

\begin{enumerate}
\item Structural properties of the coefficients
\end{enumerate}
The recursion shows inductively that

\[
\alpha_r(\lambda)\in\mathbb Q[\lambda].
\]

In fact,

\[
\deg \alpha_r=2r,
\qquad
[\lambda^{2r}]\alpha_r(\lambda)
=
\frac{3^r}{2^r r!}.
\tag{32}
\]

Indeed, in \(3\) the only contribution of degree \(2r\) comes from \(k=r-1\), through

\[
\bigl(2^2-1\bigr)\binom{\lambda-r+1}{2}
=
\frac32\lambda^2+\text{lower-degree terms}.
\]

Thus, if \(c_r\) denotes the leading coefficient of \(\alpha_r\),

\[
c_r=\frac{3}{2r}c_{r-1},
\qquad c_0=1,
\]

which gives \(32\).
Nonnegative integral parameters
If \(\lambda=m\in\mathbb N\), the expansion terminates.
Indeed, set

\[
P_m(n)=\sum_{k=0}^{m}\alpha_k(m)n^{m-k}.
\]

From \(15\) with \(N=m\),

\[
L_mP_m(n)=O(n^{-2}).
\]

But

\[
R_m(n)
=
(n+2)\bigl(P_m(n+2)-P_m(n+1)\bigr)-mP_m(n)
\]

is a polynomial in \(n\), and

\[
R_m(n)=(n+2)L_mP_m(n)=O(n^{-1}).
\]

A polynomial tending to zero along the positive integers is identically zero. Therefore \(P_m\) is an exact solution of \(1\). By uniqueness of the formal recursion,

\[
\boxed{\alpha_r(m)=0\qquad(r>m).}
\tag{33}
\]

Consequently,

\[
\lambda(\lambda-1)\cdots(\lambda-r+1)
\]

divides the polynomial \(\alpha_r(\lambda)\).

\begin{enumerate}
\item Boundary and small-parameter checks
\end{enumerate}
\(\lambda=0\)
As already noted,

\[
v_n=1,
\qquad
\alpha_r(0)=0\quad(r\ge1).
\]

\(\lambda=1\)
The formulas give

\[
\alpha_1(1)=2,
\qquad
\alpha_r(1)=0\quad(r\ge2).
\]

The corresponding exact dominant solution is

\[
v_n=n+2.
\]

Indeed,

\[
v_{n+1}+\frac{v_n}{n+2}
=
n+3+\frac{n+2}{n+2}
=
n+4=v_{n+2}.
\]

\(\lambda=2\)
Here

\[
\alpha_1(2)=7,\qquad \alpha_2(2)=10,
\qquad \alpha_r(2)=0\quad(r\ge3).
\]

Thus

\[
v_n=n^2+7n+10=(n+2)(n+5)
\]

is an exact solution. Indeed,

\[
v_{n+1}=n^2+9n+18
\]

and

\[
\frac{2v_n}{n+2}=2(n+5)=2n+10,
\]

so their sum is \(n^2+11n+28=v_{n+2}\).
\(\lambda=-1\)
The recursion gives

\[
\alpha_1(-1)=1,\qquad
\alpha_2(-1)=2,\qquad
\alpha_3(-1)=5,\qquad
\alpha_4(-1)=15.
\]

Hence

\[
v_n
=
\frac1n+\frac1{n^2}+\frac2{n^3}
+\frac5{n^4}+\frac{15}{n^5}
+O(n^{-6}).
\]

This also illustrates that negative integral values of \(\lambda\) cause no logarithmic terms in the sequence asymptotic expansion, despite the fact that logarithms can appear in the local singular form of the generating function.

\begin{enumerate}
\item Exact generating-function check
\end{enumerate}
For completeness, let

\[
U(z)=\sum_{n\ge0}u_nz^n.
\]

Multiplying \(1\) by \(n+2\), summing, and simplifying gives

\[
(1-z)U'(z)
=
(1+\lambda z)U(z)+(u_1-u_0).
\tag{34}
\]

Since

\[
H_\lambda(z)=e^{-\lambda z}(1-z)^{-\lambda-1}
\]

solves the homogeneous differential equation, every solution has, locally at \(z=0\),

\[
\boxed{
U(z)
=
e^{-\lambda z}(1-z)^{-\lambda-1}
\left[
u_0+(u_1-u_0)
\int_0^z e^{\lambda t}(1-t)^\lambda\,dt
\right].
}
\tag{35}
\]

Formula \(35\) independently confirms the dominant algebraic mode and the factorially small complementary mode. The proof above avoids having to split the singularity analysis into the nonintegral and negative-integral cases.

\begin{enumerate}
\item Reproducible coefficient generator
\end{enumerate}
The following SymPy code implements recursion \(3\). It is not used in the proof.
PythonExécuterimport sympy as sp

lam = sp.symbols("lambda")
max\_order = 8

alpha = [sp.Integer\(1\)]

for r in range\(1, max_order + 1\):
    rhs = sum(
        alpha[k]
\begin{itemize}
\item (
\end{itemize}
            (2 ** (r + 1 - k) - 1)
\begin{itemize}
\item sp.binomial(lam - k, r + 1 - k)
\item lam * (-2) ** (r - k)
\end{itemize}
        )
        for k in range\(r\)
    )

    coefficient = sp.factor(
        sp.expand\_func(sp.simplify(rhs / r))
    )
    alpha.append\(coefficient\)

for r, coefficient in enumerate\(alpha\):
    print\(f"alpha_{r} =", coefficient\)

Final conclusion
Q855 has a complete affirmative answer. For each fixed real \(\lambda\), the normalized dominant solution possesses a unique full asymptotic expansion in descending integral powers of \(n\). Its coefficients are the universal polynomials determined by \(3\), and the remainder after \(N\) terms is actually

\[
O(n^{\lambda-N-1}),
\]

rather than merely \(o(n^{\lambda-N})\).
